\begin{helper}[Support-class color state in Algorithm (A)]
Let \(\mathcal R\) be an
\noderef{ExtensionBaseRestorationTrace}, let \(q\) have exact support
\(v_n\), and suppose initially every active non-\(q\) pair is strongly
interchangeable. Assume the support-class missing-color law (A4): for the
greatest member \(t\) of every full support class \(L(w)\subseteq L^*\),
\[
 \overline\sigma_n(w)\cap\bigcup_{s\in L(w)}S_s=\{\delta_t\},
\]
and assume the active-pair overlap law (O): two exact stored pairs at active
slots \(s<t\) overlap exactly when their supports agree and no member of that
support class lies strictly between them; in the overlap case their
intersection is both \(\{\gamma_t\}\) and \(\{\delta_s\}\). Then immediately
before every step \(a\) of \(\mathcal R\),
\[
 \delta_{t_a}\in\overline{\sigma^a}(v_{t_a})
 \quad\hbox{and}\quad
 \gamma_{t_a}\in\sigma^a(v_{t_a}).
\]
\end{helper}
\begin{proof}
Fix a support \(w\) which occurs in the trace.  A trace slot supported at
\(w\) belongs to
\(L^*=\noderef{ExtensionBaseRestorationIndices}\), hence is active but is not
in the support class of \(v_n\).  By the exact chosen-vertex field in
\noderef{ExtensionBaseRestorationTrace}, this implies \(w\ne v_n\).
Conversely, every active slot supported at \(w\) is outside the support class
of \(v_n\).  Thus the full class
\(L(w)=\noderef{ExtensionBaseSupportClass}\) is a nonempty subset of \(L^*\).
Write its members, in increasing order, as
\[
                         i_0<i_1<\cdots<i_p,          \tag{1}
\]
where \(p\) is allowed to be zero.

By \noderef{ExtensionBaseActiveIndices}, every \(i_b\) is a positive slot
whose exact kind is PE.  Inspect the certificate supplied by
\noderef{TashkinovSeries}.  The kind field excludes the initial, SE, and RE
constructors of \noderef{TashkinovSeriesStep}; its PE constructor gives
\[
                    S_{i_b}=\{\gamma_{i_b},\delta_{i_b}\}.
                                                        \tag{2}
\]
Moreover \(i_b\ne q\), since otherwise the exact chosen-vertex fields would
give \(w=v_n\).  The assumed initial strong interchangeability, expanded by
\noderef{ETTStrongInterchangeable}, therefore yields
\[
                    \gamma_{i_b}\ne\delta_{i_b}.       \tag{3}
\]
For consecutive members of (1), there is no member of \(L(w)\) strictly
between them.  Applying (O), then using (2), gives
\[
              \gamma_{i_{b+1}}=\delta_{i_b}
              \qquad(0\le b<p).                       \tag{4}
\]

Apply (A4) to the greatest member \(i_p\).  It says that
\(\delta_{i_p}\) is missing at \(w\) under \(\sigma_n\), while every other
color in \(\bigcup_b S_{i_b}\) is not missing there.  The complementary
definitions \noderef{MissingColors} and \noderef{PresentColors} make each
such other color present at \(w\).  In particular, every
\(\gamma_{i_b}\) is present.  Indeed, for \(b=p\) this follows from (3).
For \(b<p\), if \(\gamma_{i_b}=\delta_{i_p}\), that color would belong to
\(S_{i_b}\cap S_{i_p}\).  Law (O) makes this intersection empty when the two
indices are not consecutive, and makes it
\(\{\delta_{i_b}\}=\{\gamma_{i_p}\}\) when they are consecutive.  The
latter alternative would give
\(\gamma_{i_b}=\delta_{i_b}\), contrary to (3).  Hence
\(\gamma_{i_b}\ne\delta_{i_p}\) in all cases.

The decreasing-order and enumeration clauses of
\noderef{ExtensionBaseRestorationTrace} show that, before the trace reaches
\(i_p\), it has processed only larger slots from other support classes.
If such a slot is \(s\), then (O) makes \(S_s\) disjoint from every
\(S_{i_b}\), because its support is not \(w\).  By
\noderef{TwoColorComponentSwap}, the switch at \(s\) changes only its two
colors.  It therefore neither changes nor creates any edge at \(w\) whose
color lies in a pair \(S_{i_b}\).  Consequently, immediately before the
step at \(i_p\), the color \(\delta_{i_p}\) is still missing at \(w\), and
every \(\gamma_{i_b}\), in particular \(\gamma_{i_p}\), is still present.

If \(p=0\), these are precisely the two assertions needed before the sole
class step \(i_0\), so this class is finished.  Suppose now that \(p>0\).
More generally, consider a descent position \(b>0\), and suppose that just
before the step at \(i_b\), the color \(\delta_{i_b}\) is missing at \(w\)
and every \(\gamma_{i_c}\) with \(c\le b\) is present there.  This stronger
induction hypothesis holds at \(b=p\) by the preceding paragraph.  Let
\(f\) be any edge incident with \(w\) and colored \(\gamma_{i_b}\).  In the graph
\noderef{TwoColorIncidenceGraph}, the edge-node of \(f\) is adjacent to the
vertex-node \(w\), since \(f\) is incident with \(w\) and its color is one
of \(\gamma_{i_b},\delta_{i_b}\).  Therefore the definition
\noderef{TwoColorComponent} places that edge-node in the connected component
rooted at \(w\).  The same argument would place every incident
\(\delta_{i_b}\)-edge in that component, but no such edge exists because
\(\delta_{i_b}\) is missing.

The trace transition at \(i_b\), together with
\noderef{TwoColorComponentSwap}, now changes every incident
\(\gamma_{i_b}\)-edge to \(\delta_{i_b}\), changes no incident
\(\delta_{i_b}\)-edge to \(\gamma_{i_b}\), and leaves all other incident
edge colors unchanged.  Hence \(\gamma_{i_b}\) is missing at \(w\) after
the switch.  By (4), used only here with \(b>0\), this color is
\(\delta_{i_{b-1}}\).

It remains to preserve the lower gamma witnesses.  For the consecutive pair,
(O) says
\[
 S_{i_{b-1}}\cap S_{i_b}
   =\{\delta_{i_{b-1}}\}=\{\gamma_{i_b}\}.
\]
Together with (3), this shows that
\(\gamma_{i_{b-1}}\notin S_{i_b}\).  Thus any incident edge witnessing the
presence of \(\gamma_{i_{b-1}}\) is fixed by the switch at \(i_b\).  If
\(c<b-1\), the indices \(i_c,i_b\) are not consecutive in their support
class, so (O) makes \(S_{i_c}\cap S_{i_b}\) empty; in particular the witness
for \(\gamma_{i_c}\) is fixed as well.  Hence every
\(\gamma_{i_c}\) with \(c<b\) remains present after the switch.

Between the steps at \(i_b\) and \(i_{b-1}\), decreasing order permits only
slots from other support classes.  Their stored pairs are disjoint from
every \(S_{i_c}\) by (O), so the same
\noderef{TwoColorComponentSwap} definition shows that they preserve both
the missing color \(\delta_{i_{b-1}}\) and all the fixed lower gamma
witnesses.  The stronger induction hypothesis, and therefore the two desired
assertions, holds just before the step at \(i_{b-1}\).  Starting at \(b=p\)
and descending proves them before every step supported at \(w\).  Since every
trace step has its stored exact support, applying this argument to that
support proves the claim for every trace index \(a\).
\end{proof}
